% SPDX-License-Identifier: CC-BY-NC-ND-4.0
% !TEX root = main.tex

\section{等化子と余等化子}

\begin{definition} \label{def: WalkingParallelPairCategory}
圏$\mathbf{P}$を$\Ob(\mathbf{P}) \coloneq 2$とし$\Mor{\mathscr{C}} \coloneq \{\id_0, \id_1, u, v\}$として定める.\
ただし$|\Mor_{\mathbf{P}}(0, 1)| = 2$である.\
これを\term{へいこうついのけん}{平行対の圏}{walking parallel pair category}という.\
平行対の圏を$\bullet \rightrightarrows \bullet$と書くこともある.
% https://q.uiver.app/#q=WzAsMixbMCwwLCIwIl0sWzEsMCwiMSJdLFswLDEsInUiLDAseyJjdXJ2ZSI6LTF9XSxbMCwwLCJcXG1hdGhybXtpZH1fMCIsMCx7ImFuZ2xlIjotOTB9XSxbMSwxLCJcXG1hdGhybXtpZH1fMSIsMCx7ImFuZ2xlIjo5MH1dLFswLDEsInYiLDIseyJjdXJ2ZSI6MX1dXQ==
\[\begin{tikzcd}
	0 & 1
	\arrow["{\mathrm{id}_0}", from=1-1, to=1-1, loop, in=145, out=215, distance=10mm]
	\arrow["u", curve={height=-6pt}, from=1-1, to=1-2]
	\arrow["v"', curve={height=6pt}, from=1-1, to=1-2]
	\arrow["{\mathrm{id}_1}", from=1-2, to=1-2, loop, in=325, out=35, distance=10mm]
\end{tikzcd}\]
\end{definition}

\begin{definition} \label{def: equalizer}
$\mathscr{C}$を圏とし$\mathbf{P}$を平行対の圏とする.\
$f, g \in \Mor_{\mathscr{C}}(A, B)$とし共変関手$E_{f,g} \colon \mathbf{P} \to \mathscr{C}$を$E_{f,g}(0) \coloneq A,\ E_{f,g}(1) \coloneq B,\ E_{f,g}(u) \coloneq f,\ E_{f,g}(v) \coloneq g$と定める.\
このときの$E_{f,g}$の極限$(\Eq(f,g), \eq(f,g))$を$f$と$g$の\term{とうかし}{等化子}{equalizer}もしくは\term{さかく}{差核}{difference kernel}と呼ぶ.\
つまり任意の$X \in \Ob(\mathscr{C})$と$\eta_0 \in \Mor_{\mathscr{C}}(X, A),\ \eta_1 \in \Mor_{\mathscr{C}}(X, B)$に対して次の図式を可換にする$\mathscr{C}$の射$\eq(f,g)_0 \backslash \eta_0 \colon X \to \Eq(f,g)$がただ1つ存在する.\
% https://q.uiver.app/#q=WzAsNCxbMSwwLCJYIl0sWzEsMiwiXFxFcShmLGcpIl0sWzAsMywiQSJdLFsyLDMsIkIiXSxbMSwyLCJcXGVxKGYsZylfMCIsMV0sWzEsMywiXFxlcShmLGcpXzEiLDFdLFsyLDMsImYiLDAseyJvZmZzZXQiOi0xfV0sWzIsMywiZyIsMix7Im9mZnNldCI6MX1dLFswLDIsIlxcZXRhXzAiLDIseyJjdXJ2ZSI6M31dLFswLDMsIlxcZXRhXzEiLDAseyJjdXJ2ZSI6LTN9XSxbMCwxLCJcXGVxKGYsZylfMCBcXGJhY2tzbGFzaCBcXGV0YV8wIiwxXV0=&macro_url=https%3A%2F%2Fgist.githubusercontent.com%2Fys-math%2F3ce50ecce6044bd5f2629781f2a56d8e%2Fraw%2Fquiver-macros.tex
\[\begin{tikzcd}
	& X & \\
	\\
	& {\Eq(f,g)} \\
	A && B
	\arrow["{\eq(f,g)_0 \backslash \eta_0}"{description}, from=1-2, to=3-2]
	\arrow["{\eta_0}"', curve={height=18pt}, from=1-2, to=4-1]
	\arrow["{\eta_1}", curve={height=-18pt}, from=1-2, to=4-3]
	\arrow["{\eq(f,g)_0}"{description}, from=3-2, to=4-1]
	\arrow["{\eq(f,g)_1}"{description}, from=3-2, to=4-3]
	\arrow["f", shift left, from=4-1, to=4-3]
	\arrow["g"', shift right, from=4-1, to=4-3]
\end{tikzcd}\]

$(\Eq(f,g), \eq(f,g)_0)$を等化子もしくは差核と呼ぶこともある.\ $\eq(f,g)_0$を$\eq(f,g)$と書くこともある.\
\end{definition}

\begin{remark} \label{rem: equalizer_is_mono}
$X \coloneq A,\ \eta_0 \coloneq \id_A$とすれば$\eq(f,g)_0$がセクションであることが分かる.\
つまり$\eq(f,g)_0$はモノ射である.\
\end{remark}

\begin{definition} \label{def: coequalizer}
$E_{f,g}$の余極限$(\Coeq(f,g), \coeq(f,g))$を$f$と$g$の\term{よとうかし}{余等化子}{coequalizer}もしくは\term{さよかく}{差余核}{difference cokernel}と呼ぶ.\
つまり任意の$X \in \Ob(\mathscr{C})$と$\eta_0 \in \Mor_{\mathscr{C}}(A, X),\ \eta_1 \in \Mor_{\mathscr{C}}(B, X)$に対して次の図式を可換にする$\mathscr{C}$の射$\eta_1 / \coeq(f,g)_1 \colon \Coeq(f,g) \to X$がただ1つ存在する.\
% https://q.uiver.app/#q=WzAsNCxbMSwwLCJYIl0sWzEsMiwiXFxDb2VxKGYsZykiXSxbMCwzLCJBIl0sWzIsMywiQiJdLFsyLDEsIlxcY29lcShmLGcpXzAiLDFdLFszLDEsIlxcY29lcShmLGcpXzEiLDFdLFsyLDMsImYiLDAseyJvZmZzZXQiOi0xfV0sWzIsMywiZyIsMix7Im9mZnNldCI6MX1dLFsyLDAsIlxcZXRhXzAiLDAseyJjdXJ2ZSI6LTN9XSxbMywwLCJcXGV0YV8xIiwyLHsiY3VydmUiOjN9XSxbMSwwLCJcXGV0YV8xIC8gXFxjb2VxKGYsZylfMSIsMV1d&macro_url=https%3A%2F%2Fgist.githubusercontent.com%2Fys-math%2F3ce50ecce6044bd5f2629781f2a56d8e%2Fraw%2Fquiver-macros.tex
\[\begin{tikzcd}
	& X & \\
	\\
	& {\Coeq(f,g)} \\
	A && B
	\arrow["{\eta_1 / \coeq(f,g)_1}"{description}, from=3-2, to=1-2]
	\arrow["{\eta_0}", curve={height=-18pt}, from=4-1, to=1-2]
	\arrow["{\coeq(f,g)_0}"{description}, from=4-1, to=3-2]
	\arrow["f", shift left, from=4-1, to=4-3]
	\arrow["g"', shift right, from=4-1, to=4-3]
	\arrow["{\eta_1}"', curve={height=18pt}, from=4-3, to=1-2]
	\arrow["{\coeq(f,g)_1}"{description}, from=4-3, to=3-2]
\end{tikzcd}\]
$(\Coeq(f,g), \coeq(f,g)_1)$を余等化子もしくは差余核と呼ぶこともある.\ $\coeq(f,g)_1$を$\coeq(f,g)$と書くこともある.\
\end{definition}

\begin{remark} \label{rem: coequalizer_is_epi}
$X \coloneq B,\ \eta_1 \coloneq \id_B$とすれば$\coeq(f,g)_1$がレトラクションであることが分かる.\
つまり$\coeq(f,g)_1$はエピ射である.\
\end{remark}
